\bookchapter{第二十八章：花對錢買到心儀機——手機選購終極指南}

\chapterintro{買手機呢件事，好多人會提前睇一週評測、比一個月價格，結果最後都係稀入麪糊塗花五六千買咗臺自己其實用唔上嘅旗艦。呢章唔列參數表——嗰啲你喺官網都能見到。我想做嘅，係畀你一套「選機思維」，以後唔管市場上出咗乜新手機，你都能自己判斷佢適唔適合你。}

\booksection{零、三步選機法：需求預算機型}

好多人喺網上問「三千蚊買乜手機好」，呢個問題其實冇得答——因為唔知你買手機主要用嚟做乜。買手機嘅正確思考順序應該係：

\begin{tcbinfo}{選機三板斧}
\begin{enumerate}
  \item \textbf{先諗清楚你最喺乎乜}：影相？打機？續航？螢幕？系統流暢？仲係淨係需要能打電話發微信就得？將需求排個序。
  \item \textbf{再定預算上限}：唔係「大概幾多錢」，係「最多能接受花幾多錢」。定一個硬上限，選機嘅時候唔會越睇越貴。
  \item \textbf{最後先至揀機型}：揸住你嘅需求排序同預算，去匹配機型。
\end{enumerate}

\end{tcbinfo}

\booksection{一、核心參數點樣睇——唔使全識，睇識呢五項就夠}

參數表動輒幾十項數據，其實普通用戶淨係需要關注呢五個維度：

\begin{center}
\begin{tabularx}{\linewidth}{|X|X|X|}
\hline
\textbf{維度} & \textbf{點樣睇} & \textbf{一句話總結} \tabularnewline
\hline
處理器（SoC） & 決定咗手機能流暢用幾年。2026年嘅標準：驍龍8 Elite Gen5/天璣9500以上=頂級；驍龍8s Gen4/天璣8400=中高端；驍龍7+ Gen4/天璣7350=夠用三年。低於驍龍6系嘅建議慎重。 & 買新唔買舊，SoC決定咗你手機能用幾年唔卡 \tabularnewline
螢幕 & OLED比LCD色彩好、更省電；更新率120Hz比60Hz掃屏流暢好多；解像度至少1080P以上。高頻PWM調光（1920Hz以上）對眼睛更友善。 & OLED+120Hz+高頻調光=好螢幕嘅及格線 \tabularnewline
鏡頭 & \textbf{唔睇像素數睇傳感器尺寸}。底大一級壓死人，傳感器尺寸比像素重要得多。超廣角實用；長焦（3x以上光學變焦）影遠處嘅嘢；微距鏡頭多數係湊數嘅。 & 大底主攝+實用長焦>多顆湊數鏡頭 \tabularnewline
電池同充電 & 4500mAh係基礎及格線，5000mAh以上一日一充好穩。快充睇瓦數：60W以上算快，100W以上飛快。無線充電方便但慢。 & 5000mAh+60W快充=唔焦慮組合 \tabularnewline
系統更新承諾 & 決定咗手機能安全用幾耐。蘋果一般畀5-6年系統更新；三星旗艦畀4年系統+5年安全更新；谷歌Pixel畀7年；國產中端機更新承諾普遍偏短（2-3年）。 & 想用耐啲，選系統更新承諾長嘅品牌 \tabularnewline
\hline
\end{tabularx}
\end{center}

\booksection{二、按預算分檔推薦（以2026年中市場為基準）}

以下推薦基於2026年中國內地市場行情，價格為官方起售價。具體價格會有浮動，以你入手時嘅實際價格為準。

\booksubsection{1000元以下：能打電話，能掃微信，夠用就得}

\textbf{適合邊個}：學生、後備機、對手機要求唔高、主要用嚟打電話同微信。

\textbf{能買到乜}：Redmi Note系列、榮耀X系列、realme數字系列嘅中低配。處理器大概驍龍6系或天璣6000系列，螢幕係LCD為主，影相能影清楚但對焦慢，系統流暢度一般。

\textbf{選購建議}：呢個價位睇兩點：電池夠唔夠大（選5000mAh以上嘅），記憶體夠唔夠（選6GB+128GB起步，唔好買4GB嘅會卡）。紅米喺呢個價位優勢最明顯。

\booksubsection{1000-2000元：性價比最高，能用三四五年唔卡嘅甜點區}

\textbf{適合邊個}：學生、預算有限但想用得耐、對影相同遊戲冇極致追求。

\textbf{代表性機型}：Redmi Turbo系列、iQOO Z系列、榮耀X系列高配、一加Ace系列入門。

\textbf{選購建議}：選搭載驍龍7+ Gen3或天璣7300以上處理器嘅機型，配8GB+256GB記憶體。呢個配置日常使用+偶爾打王者榮耀完全夠用，正常能用3-4年唔卡。

\booksubsection{2000-3500元：中堅力量，效能同影相都有明顯提升}

\textbf{適合邊個}：想要唔錯嘅影相、打機唔卡、希望用得耐嘅普通用戶。呢個係最多人買嘅價格區間。

\textbf{代表性機型}：小米數字系列（如小米17標準版降價後）、vivo S系列高配、OPPO Reno系列、榮耀數字系列、Redmi K系列。

\textbf{選購建議}：呢個價位最大嘅特點係「水桶機」——冇明顯短板。比起2000元以下嘅機仔，主要提升喺：相機明顯變好（有光學防震喇）、螢幕上咗高素質OLED、快充更快（80W以上常見）、做工質感亦更精緻。選12GB+256GB或12GB+512GB，用四五六年冇問題。

\booksubsection{3500-5000元：準旗艦，除咗極客玩家其他人完全夠用}

\textbf{適合邊個}：對影相有追求嘅用戶、想體驗頂級質感嘅用戶。

\textbf{代表性機型}：小米數字系列標準版、vivo X系列、OPPO Find X系列標準版、iPhone舊款旗艦（如iPhone 16）。

\textbf{選購建議}：呢個價位已經能買到舊年嘅旗艦機型。我嘅建議係：\textbf{買舊年嘅旗艦，唔買今年嘅中端}。舊年旗艦嘅用料、做工、鏡頭硬件通常比今年嘅中端高一個檔次。前提係你接受佢唔係「最新款」。

\booksubsection{5000元以上：旗艦區，唔差錢或對某項體驗有極致追求}

\textbf{適合邊個}：攝影發燒友、重度遊戲玩家、追求頂級螢幕同工藝嘅用戶、用手機影影片做內容嘅創作者。

\textbf{代表性機型}：iPhone 17 Pro/Pro Max、三星S26 Ultra、華為Pura 80 Ultra、小米17 Ultra、vivo X300 Pro。

\textbf{選購建議}：呢個價位各家嘅差距越嚟越細，其實主要睇你鍾意乜系統。鍾意影相選vivo華為，鍾意影影片選iPhone，鍾意螢幕同SPen選三星，鍾意均衡選小米。講真到咗呢個價位，\textbf{你可以任性，選你睇着順眼嘅嗰部就得}。

\booksection{三、按場景選機——你嘅需求係乜？}

以下係幾個最常見嘅「使用畫像」，對號入座就得：

\begin{center}
\begin{tabularx}{\linewidth}{|X|X|X|}
\hline
\textbf{使用場景} & \textbf{重點睇乜} & \textbf{推薦方向} \tabularnewline
\hline
影相達人 & 大底主攝（1吋最佳）、光學防震、長焦（3x以上光學變焦）、影像晶片 & vivo X系列、華為Pura/Mate系列、小米Ultra系列 \tabularnewline
遊戲玩家 & 頂級處理器、高觸控採樣率、散熱系統、大電池+快充 & iQOO數字系列、ROG Phone、紅魔、Redmi K系列 \tabularnewline
續航優先 & 5500mAh以上電池、能效好嘅處理器、LTPO螢幕 & 榮耀Magic系列、一加、Redmi Note大電池版 \tabularnewline
長輩使用 & 簡易模式、大字體、大音量、大電池、售後網點多 & 華為暢享、榮耀Play、OPPO A系列、Redmi Note \tabularnewline
學生用機 & 性價比高、耐用、電池大、有家長控制功能 & Redmi Note/Turbo、realme GT、iQOO Neo \tabularnewline
蘋果生態用戶 & 已有Mac/iPad/AirPods，買iPhone係唯一解 & iPhone 17標準版（性價比）或Pro版（滿血體驗） \tabularnewline
出國旅行 & eSIM支援、全球頻段覆蓋完整、Google服務 & iPhone、三星S系列、Pixel、一加 \tabularnewline
\hline
\end{tabularx}
\end{center}

\booksection{四、五個容易踩嘅坑——買之前對照睇一眼}

\begin{enumerate}
  \item \textbf{坑一：四攝五攝唔等於影相好}——廠商鍾意喺參數上堆鏡頭數量，入面成日有兩顆係200萬像素嘅「湊數鏡頭」（微距+景深），實際完全冇用。睇手機影相能力，\textbf{淨係睇主攝+超廣角+長焦}呢三顆就夠。其他多出嚟嘅可以當佢唔存在。
  \item \textbf{坑二：唔好被1億像素唬住}——像素高唔代表畫質好。傳感器尺寸、光圈、算法調教比像素重要得多。好多旗艦機用5000萬像素嘅大底傳感器，影出嚟嘅相比某啲1億像素嘅千元機好太多。
  \item \textbf{坑三：唔好買運營商定製機}——「存話費送手機」睇落好劃算，但運營商定製機系統更新慢、可能有刪唔到嘅預裝App、而且通常係運營商鎖定嘅（淨係用得呢間運營商嘅卡）。睇似慳咗錢，實則「鎖定」咗你。
  \item \textbf{坑四：「曲面屏好高級」——係營銷，唔係真理}——曲面屏曾經係旗艦機嘅標緻，但而家越來越多人發現佢嘅問題：貼膜難、邊緣誤觸、爛咗維修貴。2026年嘅趨勢係回歸「微曲」甚至「純直屏」。唔好因為「曲面嘅就係高級嘅」而選曲面屏。
  \item \textbf{坑五：唔係所有「AI手機」都真係有用}——2026年所有廠商都喺度喊「AI手機」，但好多淨係將已有嘅語音助手換咗個名。真正有用嘅AI功能係：實時翻譯、AI消除（影相將人去掉）、AI修圖、端側大模型推理、本地化嘅智能摘要。其他嘅花哨功能基本都係噱頭。
\end{enumerate}

\booksection{五、幾時買最劃算？}

\begin{itemize}
  \item \textbf{618（5月底到6月中）}：一年中手機降價幅度最大嘅時段。好多手機廠商為咗衝銷量，旗艦機能降價800-1500元。係上半年買手機嘅最佳視窗。
  \item \textbf{雙11（10月底到11月中）}：下半年嘅大促，力度同618差唔多。呢個時間段好多品牌新一代旗艦啱啱發布，上一代旗艦降價幅度最大。
  \item \textbf{新機發布後一個月}：啱啱發布時原價，一個月後渠道價（第三方舖頭）一般會平200-500。
  \item \textbf{春節前兩週到元宵節}：線下實體舖清庫存嘅「年貨節」，可以同店員傾價。
  \item \textbf{唔推薦嘅時間}：新品發布會後當日衝首發——除非你真係特別想做第一批用戶，否則首發價係最貴嘅。
\end{itemize}

\booksection{六、渠道點樣選——喺邊買靠譜？}

\begin{itemize}
  \item \textbf{品牌官方商城/天貓旗艦店/京東自營}：最保險，價格冇優勢但唔使擔心翻新同假貨，售後最好。
  \item \textbf{拼多多「百億補貼」}：價格最低，但需要選標註咗「百億補貼」嘅舖頭。冇標註嘅唔好碰。收到貨記住錄開箱影片，有問題能退貨。
  \item \textbf{線下實體舖}：可以親手摸、試手感。注意：導購會推利潤高嘅機型（而唔係最適合你嘅），建議入舖之前先喺網上鎖定2-3款目標機型，去到就話「我淨係想睇呢幾款」。
  \item \textbf{二手平臺（閒魚/轉轉）}：平但有風險。如果要買二手，記住三條：要求提供購買發票或訂單截圖、要求提供IMEI號查驗、\textbf{淨係走平臺交易，絕唔私下轉賬}。
\end{itemize}

\begin{tcbpractice}{動手：用呢套方法選一部手機}
\begin{enumerate}
  \item 攞一張紙，寫低你嘅兩個核心需求（比如：續航第一、影相第二）。
  \item 定一個預算上限（比如3500元）。
  \item 打開京東或天貓，按價格篩選，圈出3-5款備選機型。
  \item 去B站搜「XX機型 一個月使用體驗」（唔係第一日評測，係用咗至少一個月之後嘅真實感受），睇兩三條片。
  \item 確認渠道靠譜後落單，收貨時錄開箱影片。
\end{enumerate}

\end{tcbpractice}

\begin{tcbsummary}{本章小結}
\begin{itemize}
  \item 選機三板斧：先諗需求 再定預算 最後揀機型
  \item 核心參數淨係睇五項：處理器、螢幕、鏡頭、電池充電、系統更新承諾
  \item 1500-3500元係性價比最高、最推薦嘅甜點區間
  \item 買舊年嘅旗艦比買今年嘅中端更劃算
  \item 五大坑：多攝騙局、高像素陷阱、運營商定製、曲面屏迷信、AI噱頭
  \item 618同雙11係一年中最適合買手機嘅兩個時間視窗
\end{itemize}

\end{tcbsummary}
